\section{The Basic Finite Construction}
\label{sec:basic-construction}

We illustrate the central idea for one evaluation of a public Boolean circuit
$F$.  The input providers are named application parties and the output is
public.  Computation roles are assigned to registered computation keys by a
public rule.  The construction uses the ledger only through the generic API
of \cref{sec:model}; garbling, role assignment, encryption, slot closure, and
message validation are protocol operations.  Our first proof target fixes the
corrupt set of named input providers before setup.  The adaptive interface of
\cref{sec:model} is retained as a later target.

This section is a bounded, two-phase blueprint.  It cleanly isolates the use
of silence, but does not by itself instantiate the sampler that correlates its
pre-input and post-input phases.  That requirement becomes property~$Z$ in
\cref{sec:proof-outline}; the streaming candidate of \cref{sec:ongoing}
implements the same schedule repeatedly.

\subsection{Public structural branches}

Fix a ledger prefix $L_0$ containing the registered input providers and at
least $2mk$ fresh computation keys, where $m$ is the number of input wires.
For each wire $w$, the public role map selects two disjoint size-$k$
committees $C^w_0$ and $C^w_1$.  All $2m$ committees are mutually disjoint,
and no computation key is assigned another canonical role.  Their subscripts
are \emph{structural branch indices}; they do not denote semantic bits.

For concreteness, the generic setup distribution may place a finite sampler
parameter $U_{\mathsf{fin}}$ in $\mathsf{crs}$.  Before inputs arrive, its
preparation phase samples for each input wire
\[
  (K^w_0,K^w_1),\qquad (T^w_0,T^w_1),\qquad
  \pi_w\sample S_2 .
\]
The garbling uses the semantic assignment
\[
  b\longmapsto K^w_{\pi_w(b)}.
  \label{eq:semantic-assignment}
\]
Equivalently, one may view it as garbling the structurally reindexed circuit
\[
  F_{\boldsymbol\pi}(z)=
  F(\pi_1^{-1}(z_1),\ldots,\pi_m^{-1}(z_m)).
\]
The pre-input public object $\mathsf{Prep}_F$ contains
the encrypted role packets described below, but not yet the garbled tables.
The complete label pairs and semantic maps remain hidden in the finite
sampler's logical state.

For public token validation, preparation publishes the branch-indexed tags
\[
  V^w_s=h_{\mathsf{tok}}(\mathsf{sid},w,s,T^w_s),
  \qquad s\in\{0,1\}.
  \label{eq:token-tags}
\]
Tokens have high min-entropy, and $h_{\mathsf{tok}}$ is assumed one way and
collision resistant on the domain-separated token inputs.  Thus the tags
allow everyone to recognize structural branch $s$ without revealing the
token in advance or the semantic value of that branch.

Label $K^w_s$ is verifiably $(t+1)$-out-of-$k$ shared among $C^w_s$, where
$t=\lfloor(k-1)/3\rfloor$.  For each member with public key $pk_i\in C^w_s$,
$\mathsf{Prep}_F$ contains
\[
  D_i=\mathsf{Enc}_{pk_i}
      (w,s,T^w_s,\mathsf{share}_{i,s},
       \mathsf{verification\ data}).
  \label{eq:finite-role-packet}
\]
It similarly contains, under the registered key of the application party
providing wire $w$, an encrypted packet specifying the map
$b\mapsto T^w_{\pi_w(b)}$.  Packet lengths, ordering, verification data, and
ciphertext distributions are symmetric in the two structural branches.
Neither $\pi_w$ nor the semantic bit represented by a public committee is
revealed.

The fixed $\mathsf{WhosTurn}$ waits until $L_0$ exists and then schedules the
provider keys and role keys in a canonical order.  Registration itself has no
role argument.

\subsection{Input activation and irrevocable default}

When the provider for wire $w$ is scheduled, an honest provider with input
$x_w$ supplies a protected algorithm $A_{w,x_w}$ to the ledger.  The
algorithm derives $\mathsf{Prep}_F$, decrypts the provider packet using
$sk_P$, and returns the tagged structural token $T^w_{\pi_w(x_w)}$.  The
ledger appends this output without revealing $sk_P$, the other token, or the
input-dependent algorithm.

A corrupt provider may instead put any string in its turn.  At a fixed cutoff
the protocol derives one irrevocable close record $m_w$.  Define
\[
  \mathsf{Valid}_w(m)=s
  \quad\Longleftrightarrow\quad
  h_{\mathsf{tok}}(\mathsf{sid},w,s,m)=V^w_s
\]
for exactly one $s\in\{0,1\}$; otherwise set
$\mathsf{Valid}_w(m)=\bot$.  A $\mathsf{NoMsg}$ record, malformed string,
ambiguous match, and late entry all give $\bot$.  The effective structural
branch is
\[
  s_w=\begin{cases}
       \mathsf{Valid}_w(m_w),&\mathsf{Valid}_w(m_w)\neq\bot,\\
       0,&\mathsf{Valid}_w(m_w)=\bot.
      \end{cases}
  \label{eq:finite-default}
\]
Because $\pi_w$ is uniform and hidden, default branch $0$ has semantic value
$\pi_w^{-1}(0)$, a uniform bit.  A corrupt provider that has opened its own
packet knows the semantic map and may choose its effective input, just as a
corrupt ideal input party may.

The roles, not a protocol-aware ledger, implement this rule.  Committee
$C^w_s$ speaks precisely if $\mathsf{Valid}_w(m_w)=s$, or if
$\mathsf{Valid}_w(m_w)=\bot$ and $s=0$.  In particular, $C^w_0$ can
distinguish a valid branch-$1$ token from a malformed token using the public
tags and therefore does not speak in the former case.  Closure must precede
the default release: otherwise a malformed message followed by a late valid
token could expose both labels.

\subsection{Conditional key release and evaluation}

Each computation member is scheduled once.  Its honest protected algorithm
derives its canonical $D_i$, decrypts it with $sk_i$, checks the public close
record and token tags, and returns $sk_i$ exactly under the rule above;
otherwise it returns $\bot$.  The selected committee therefore reveals its
keys, while every honest member of the sibling committee produces only a
$\mathsf{NoMsg}$ timeout and keeps its key hidden forever.

The ledger performs no validation.  Protocol observers ignore a purported
key unless $\mathsf{KeyMatch}(pk_i,sk_i)$ holds, decrypt $D_i$ themselves,
and reject malformed shares using the public verification data.  Conditioned
on $\mathsf{Good}$, the selected committee has enough honest members to reveal
at least $t+1$ valid shares despite corrupt withholding.  Everyone
reconstructs exactly one active structural label: $K^w_{\pi_w(x_w)}$ for a
valid honest input and $K^w_0$ for a defaulted input.

Only after all input slots close does the completion phase emit the public
garbled tables $\widetilde F$ and output-decoding information, using the
label pairs and assignment in \eqref{eq:semantic-assignment}.  All parties
evaluate $\widetilde F$ from the reconstructed active labels and decode the
public output.  This order is important for the intended straight-line
simulation: the ideal public output is available before the tables must be
fixed.  Publishing the tables before input closure would instead require an
online or adaptively simulatable garbling primitive.

Ordinary garbling decoding is sufficient.  A protocol needing a private bit
$y$ for an honest recipient can, at the application layer, let the recipient
input a fresh pad $r$ and publicly output $y\oplus r$; version~1 makes no
claim for later adaptive corruption of that pad holder.

\subsection{Why complete exposure need not be fatal}

Publishing $sk_i$ exposes the selected role's complete retained state,
including its packet plaintext, local coins, and call record.  The exposed
share belongs to a label reconstructed publicly anyway and gives no share of
the disjoint sibling label.  In the sibling committee, at most $t$ members are
corrupt on $\mathsf{Good}$, while every honest member remains silent;
threshold privacy therefore protects the unselected label.

The hidden permutation is necessary in this public-role model.  Without it,
the speaking committee would reveal the input.  With it, the public structural
branch has no fixed semantic meaning, and garbling privacy handles the
remaining dependence on the input.

The static-corruption restriction on named providers is used here.  A
provider corrupt before setup controls its slot; an honest provider is never
later asked to explain the encrypted unused token or its hidden semantic map.
An adaptive extension will require an equivocal or non-committing treatment of
the provider packet.

This remains a blueprint rather than a theorem.  A proof must combine the
two-stage finite-sampler property, multi-user encryption, token-tag security,
verifiable-sharing privacy and robustness, garbling privacy, and the
good-committee bound.  It must also account for retrospective decryption.
One-role-per-key gives one canonical packet, not one ciphertext globally:
adversarial evaluations of a copyable sampler may address further fork
packets to the same $pk_i$.  Canonical-root isolation must ensure that every
such plaintext is independent fork material containing no canonical inactive
label; the readout part of property~$Z$ must simulate all openings.

\paragraph{A gate-by-gate variant.}
For intuition, garbling can be replaced by an expensive
information-theoretic construction.  For every NAND gate and public
structural branch corresponding to a pair of input tokens, a separate
committee holds shares of the corresponding output token.  Only the committee
whose branch matches the public inputs reveals its keys.  This makes the role
of silence especially transparent, at the cost of a large setup and many
registered computation parties.
